% Provenance: consolidation plan v3, advisor brief, review verdicts.
\chapter{Conclusions, Limitations, and Future Work}
\label{ch:conclusion}

\section{What this thesis established}
Under binary verifiable rewards, group-relative training starves at both ends
of the difficulty spectrum, and the Zero-Variance Fraction --- read jointly
with mean reward --- is a near-free, calibrated instrument that sees both
walls, including the one reward curves cannot show (Claim~1). Group size
selects which wall a fixed rollout budget hits: small $G$ buys early optimiser
steps and starves in the endgame; large $G$ holds signal throughout
(Claim~2). The estimation theory supporting the instrument is calibrated
(Wilson CIs), its waiting-time budget is empirically exact at ratio 1.00, and
its would-be optimal-group-size theory resolves to an honest negative result:
per-rollout accounting always answers $G^\star\in\{2,3\}$, for any data, so
adaptivity must be justified by richer objectives.

\section{Limitations}
All training-side evidence sits on one model (Qwen3-8B, LoRA rank 4), one
task family (GSM8K, binary rewards), one managed API with a closed loss
kernel, at 1--3 seeds per result; pre-2026-07-11 artifacts additionally carry
a since-fixed off-by-one target and v1 parser semantics. \ZVF{} as defined is
specific to (near-)discrete rewards; continuous-reward analogues (e.g.\
contrast-density statistics) are sketched but not validated. The pooled
cross-scale correlations are descriptive only. Theory carries declared proof
gaps. None of the seventeen working papers is submission-ready as a standalone
unit --- by design, after external review (below).

\section{Publication roadmap (post-degree, gated)}
Following three rounds of external review, the programme's publication plan
is thesis-first: (i) a bounded diagnostic paper (Claims~1--2, this thesis's
Chapters~4--6 core) is publishable without any controller claim; (ii) the
controller paper is gated on a pre-registered, compute-matched win over
static $G{=}16$ \emph{and} a reward-only schedule at $\ge3$ seeds;
(iii) the survival audit is gated on an open stack and differential-objective
test tooling; (iv) the benchmark remains a versioned artifact. The
seventeen-document corpus maps onto these vehicles and this thesis per
Appendix~A.

\section{Future work}
Beyond the gates above: MATH-500 and code-task completion of the capability
panels; a clustered-variance CI for curriculum-ordered training; a
variance-aware replacement for the T3 objective; continuous-reward \ZVF{}
analogues; and the full bridge-to-DPO analysis of the group-size working
paper.

\section{Closing}
The programme's most transferable lesson is methodological: at this scale,
the binding constraint on RL post-training research is not compute or
algorithmic novelty but \emph{certainty about what actually ran}. A cheap
diagnostic that cannot be gamed by a pretty reward curve, a reporting standard
whose every item once flipped a result, and the habit of reading the runner
--- these travelled further than any single number in this thesis.
